\documentclass[11pt,a4paper]{article}

\usepackage[margin=1.05in]{geometry}
\usepackage{amsmath,amssymb,amsthm}
\usepackage{mathtools}
\usepackage{amscd}
\usepackage{booktabs}
\usepackage{listings}
\usepackage{xcolor}
\usepackage[colorlinks=true,linkcolor=black,citecolor=black,urlcolor=blue]{hyperref}
\usepackage{enumitem}
\usepackage[T1]{fontenc}
\usepackage[protrusion=true,expansion=false]{microtype}
\setlength{\emergencystretch}{2.5em}
% msbm has no double-struck digits, so \Id renders as a wrong glyph;
% this is the identity operator throughout.
\newcommand{\Id}{\mathrm{1}\hspace{-0.30em}\mathrm{l}}

%%%%%%%%%%%%%%%%%%%%%%%%%%%%%%%%%%%%%%%%%%%%%%%%%%%%%%%%%%%%%%%%%%%%%%%%%%%%%%%
% Theorem environments
%%%%%%%%%%%%%%%%%%%%%%%%%%%%%%%%%%%%%%%%%%%%%%%%%%%%%%%%%%%%%%%%%%%%%%%%%%%%%%%
\theoremstyle{plain}
\newtheorem{theorem}{Theorem}[section]
\newtheorem{proposition}[theorem]{Proposition}
\newtheorem{lemma}[theorem]{Lemma}
\newtheorem{corollary}[theorem]{Corollary}
\newtheorem{conjecture}[theorem]{Conjecture}

\theoremstyle{definition}
\newtheorem{definition}[theorem]{Definition}
\newtheorem{example}[theorem]{Example}
\newtheorem{principle}[theorem]{Principle}
\newtheorem{obligation}[theorem]{Obligation}

\theoremstyle{remark}
\newtheorem{remark}[theorem]{Remark}

%%%%%%%%%%%%%%%%%%%%%%%%%%%%%%%%%%%%%%%%%%%%%%%%%%%%%%%%%%%%%%%%%%%%%%%%%%%%%%%
% Macros.  NOTATION RULES OBSERVED THROUGHOUT:
%   * the rho calculus is never written with the Greek letter rho
%   * quoting is @P, dropping is *x; corner quotes are pre-2016 notation
%%%%%%%%%%%%%%%%%%%%%%%%%%%%%%%%%%%%%%%%%%%%%%%%%%%%%%%%%%%%%%%%%%%%%%%%%%%%%%%
\newcommand{\rhoc}{\textsf{rho}}
\newcommand{\grho}{\textsf{rho}_{\V}}
\newcommand{\sem}[1]{[\![#1]\!]}
\newcommand{\V}{\mathbb{V}}
\newcommand{\Bool}{\mathbb{B}}
\newcommand{\Rnn}{\mathbb{R}_{\geq 0}}
\newcommand{\Cx}{\mathbb{C}}
\newcommand{\Nat}{\mathbb{N}}
\newcommand{\Proc}{\mathcal{P}}
\newcommand{\Names}{\mathcal{N}}
\newcommand{\NF}{\mathrm{NF}}
\newcommand{\Reach}{\mathrm{Reach}}
\newcommand{\supp}{\mathrm{supp}}
\newcommand{\Per}{\mathrm{Per}}
\newcommand{\Det}{\mathrm{Det}}
\newcommand{\scong}{\equiv}
\newcommand{\tensorv}{\otimes}
\newcommand{\sumv}{\oplus}
\newcommand{\Cand}{\mathrm{Cand}}
\newcommand{\Out}{\mathrm{out}}
\newcommand{\crisp}[1]{\lceil #1 \rceil}
\newcommand{\pow}{\mathrm{pow}}
\newcommand{\str}{\mathrm{s}}
\newcommand{\conf}{\mathrm{c}}
\newcommand{\ket}[1]{\lvert #1 \rangle}

\lstdefinestyle{rholang}{
  basicstyle=\ttfamily\small,
  keywordstyle=\bfseries,
  morekeywords={for,new,in,where,contract,match,if,else},
  columns=fullflexible,
  keepspaces=true,
  showstringspaces=false,
  breaklines=true,
  frame=single,
  framesep=5pt,
  rulecolor=\color{black!35},
  xleftmargin=6pt,
  aboveskip=9pt,
  belowskip=9pt
}
\lstset{style=rholang}

%%%%%%%%%%%%%%%%%%%%%%%%%%%%%%%%%%%%%%%%%%%%%%%%%%%%%%%%%%%%%%%%%%%%%%%%%%%%%%%

\title{\bfseries Races Decided by the Born Rule\\[4pt]
\large The complex instance of graded where-clauses: a quantum namespace,\\
contention-set completion, and the fragment on which a process term\\
denotes a quantum state}

\author{L.\,G.\ Meredith\\ \small F1R3FLY.io}

\date{F1R3FLY.io working note --- August 2026\\[2pt]
\normalsize Version 5. Paper I of three; supersedes the quantum half of
\emph{Graded Where-Clauses}}

\begin{document}
\maketitle

\begin{abstract}
The rho calculus does not say how the non-determinism of a race is resolved. An
implementation must supply that mechanism, and quantum mechanics is one of the
available choices: it supplies amplitudes for the branches of a race and the
Born rule for turning them into outcomes. This note takes that instantiation
seriously and asks what has to be true of a program before its denotation is a
quantum state.

The design is one line of grammar, carried over from the companion paper on
resolution algebras: the \texttt{where} clause of a for-comprehension is valued
in a commutative semiring rather than in the Booleans, and is evaluated per
candidate match. Here $\V = \Cx$ throughout.

Four things are new. First, a \emph{quantum namespace}: names bound by
\texttt{new} are unforgeable and are the only ones on which a clause may take
complex values, while quoted names are classical. Unique ownership of quantum
data --- which existing quantum process calculi obtain from a type system ---
follows from unforgeability, and the herald of a non-unitary step lands in the
classical namespace by construction, which is what gives measurement a place to
happen. Second, the denotation is \emph{register-indexed}: $\sem{P}$ lives in a
Fock space over the names $P$ can act on, and parallel composition combines
states by adding occupancies. Because parallel composition is
associative-commutative, the multi-occupancy sectors are symmetric, so bosonic
statistics and the permanent of the scattering matrix are forced by structural
congruence rather than chosen. Third, the completion condition that makes a step
norm-preserving is stated over the \emph{contention set} --- the connected
components of the conflict relation, with maximal matchings as the alternatives
and the fibre-summed amplitudes normalised --- and conflict components tensor
rather than sum. The race equation, the join, and the two-receipt race are the
special cases. Fourth, we identify the fragment on which the one-step operator
is an isometry. It is not the whole calculus: absorbing normal forms and
spectator tokens each break it, and both are visible syntactically. Under halt
tagging, staging, and a dual-rail encoding the failures disappear in every case
we have been able to construct.

The dual-rail encoding also settles a question the previous version got wrong.
Under the payload encoding, a step that uses each bound name exactly once is
block-diagonal in the input basis, hence not a gate; that is a fact about the
encoding, and under dual rail every $2\times 2$ and the entangling
$4 \times 4$s are realised with no discard anywhere. Shor's algorithm is
rebuilt on that basis, with the whole term given, generated from the same gate
list the simulator runs.

Non-unitary clauses are contained rather than excluded: contractivity is a
judgement, the defect is dilated at the site, and the ancilla is neither traced
nor projected but written as a classical message which a retry loop may read
and the resource ledger charges. What was free post-selection becomes a bill of
$2^n$ firings.

We are explicit about what is not established: norm preservation is proved at a
configuration and conjectured for arbitrary concurrency; the staging condition
is a measured survey and not a theorem; the complexity containment is a
conjecture; and quantum alternation is known to be incompatible with the
standard fixed-point semantics for recursion, which our budget does not repair.
\end{abstract}

\tableofcontents

%%%%%%%%%%%%%%%%%%%%%%%%%%%%%%%%%%%%%%%%%%%%%%%%%%%%%%%%%%%%%%%%%%%%%%%%%%%%%%%
\section{Introduction}\label{sec:intro}

\subsection{Where the choice is}

A rho program does not determine its own execution. Write
\begin{equation}\label{eq:race}
  x!(Q_1) \mid \textsf{for}(y \leftarrow x)P \mid x!(Q_2)
\end{equation}
and the operational semantics permits two reductions. It does not say which one
happens, and it does not say with what tendency. This is not an oversight; it is
the defining feature of a mobile process calculus, and it is what makes the
calculus a specification of interaction rather than of a machine.

An implementation must close the gap. RSpace closes it by presenting the
parallel fragment as a store from channels to polarity-homogeneous bags and
letting the physical race between cores decide which produce is zipped with
which consume. That is a resolution mechanism, and it is a choice.

The resolution mechanism is a \emph{parameter}, and other values of the
parameter are available. Quantum mechanics is one: it supplies amplitudes for
the branches of a race and the Born rule for turning them into outcomes. This
note develops that instantiation. The companion paper \cite{ResAlg} develops the
parameter itself --- the resolution algebra, the conservativity of the Boolean
value, the normal-form theorem for weight maps --- and a second companion
\cite{PLNPaper} develops the non-negative real instance and its use in
inference-driven execution. Everything here is $\V = \Cx$.

\subsection{What this paper must answer}

The previous version of this material \cite{GWv4} was reviewed, and the review
was right about four things. Non-contractive clauses were legal but had no
dilation. Dilating a receipt site does not make the one-step operator
norm-preserving, because the operator sums over sites as well as over
candidates. The herald was simultaneously a coherent ancilla and an ordinary
message, with no account of the transition between them. And the stage-one
proposition had a counterexample.

Two further defects were not raised and are addressed here. Norm preservation of
each column is only half of $T^{\dagger}T = \Id$; the other half is that
distinct columns be orthogonal, and that fails for reasons having nothing to do
with weights. And a process term is not, in general, the carrier of a quantum
state at all --- states must be indexed by a register, and states over
overlapping registers do not combine.

\subsection{Contributions}

\begin{enumerate}[leftmargin=1.4em,itemsep=3pt]
\item \textbf{The quantum namespace} (\S\ref{sec:namespace}). Names bound by
\texttt{new} are unforgeable; quoted names are not. Take the first as the
quantum namespace. Unique ownership follows from unforgeability rather than from
a type discipline, and the herald channel $h(r) := @r$ is classical because it
is a quotation --- so the classical/quantum boundary, which is where measurement
happens, is drawn by the calculus and not by us.

\item \textbf{The denotation is register-indexed} (\S\ref{sec:register}).
$\sem{P}$ lives in a Fock space over the names $P$ can act on. Composition adds
occupancies, and is injective exactly when the registers are disjoint;
overlapping registers forget which side supplied which quanta, which is
indistinguishability. Since $\mid$ is associative-commutative the sectors are
symmetric, so \textbf{bosonic statistics are forced by structural congruence}
(Proposition~\ref{prop:bose}).

\item \textbf{Contention-set completion} (\S\ref{sec:completion}). The unit of
normalisation is the connected component of the conflict relation; the
alternatives within a component are its maximal matchings; the fibre-summed
amplitudes are what must have unit norm; and components tensor rather than sum.
Proposition~\ref{prop:complete} and Table~\ref{tab:conventions}, which measures
the three wrong answers against the right one.

\item \textbf{The quantum-safe fragment} (\S\ref{sec:fragment}). Column
orthogonality fails in two ways, both syntactic: \emph{absorption}, where a
normal form idles into a configuration another branch is still flowing towards,
and \emph{spectators}, where a parked token is indistinguishable from an emitted
one. Halt tagging removes the first, staging removes the second, and dual rail
removes both. Table~\ref{tab:collisions}.

\item \textbf{Linearity does not cost universality} (\S\ref{sec:gadgets}). The
previous version proved that a step using every bound name exactly once is
block-diagonal in the input basis. The proposition is true and its hypothesis is
a property of the encoding. Under dual rail every $2\times 2$ is realised with
no discard (Lemma~\ref{lem:1q}), and so is every entangling $4\times4$
(Lemma~\ref{lem:2q}).

\item \textbf{Shor, in full} (\S\ref{sec:shor}). The whole term, gadget
contracts and instance, generated from the gate list the simulator runs.

\item \textbf{Heralds and the price} (\S\ref{sec:dilation}). Contractivity is a
judgement; the defect is dilated at the site; the ancilla is written as a
classical message; post-selection becomes a retry loop that the ledger charges
at $2^n$ firings.

\item \textbf{Position} (\S\ref{sec:related}). The right comparison class is
quantum control, not quantum process calculus. Two of Bădescu and Panangaden's
negative results \cite{BadescuPanangaden} are discharged by the design and one
is not; their suggested way out of the quantum-alternation impasse ---
interferometry with no distinct control qubit --- is this design.
\end{enumerate}

\subsection{What this paper does not do}

It does not prove a quantum advantage. It does not prove
Conjecture~\ref{conj:ledger}. It proves norm preservation at a configuration and
conjectures it for arbitrary concurrency (Conjecture~\ref{conj:global}). The
staging condition of \S\ref{sec:fragment} is a measured survey over small
systems, not a theorem. And it does not solve the recursion problem of
Obligation~\ref{ob:recursion}.

%%%%%%%%%%%%%%%%%%%%%%%%%%%%%%%%%%%%%%%%%%%%%%%%%%%%%%%%%%%%%%%%%%%%%%%%%%%%%%%
\section{Why context-based full abstraction cannot help}\label{sec:fa}

Going from the pi calculus into rho, the question of how non-determinism is
resolved never arises, because both languages are silent in the same way.
Mapping rho to quantum mechanics is different: quantum mechanics makes a
definite commitment about resolution, namely the Born rule. A repair that
enlarges the class of rho contexts therefore cannot expose the additional
structure, because contexts can \emph{create} apertures --- add contention ---
but cannot \emph{read} them. Every context is itself resolved by the same
unspecified mechanism.

This is not merely an observation about rho.

\begin{proposition}[No extensional semantics, after
\cite{BadescuPanangaden}]\label{prop:noext}
A conditional that applies a global phase in one branch and the identity in the
other implements a controlled phase, although the two branches are physically
indistinguishable as quantum operations. Hence there is no structural semantics
for quantum alternation based on operations with extensional equality.
\end{proposition}

Our situation is the process-calculus shadow of that statement: the resolver is
exactly the datum that an extensional semantics discards, so no amount of
context can recover it.

\begin{definition}[Resolver separation]\label{def:sep}
Fix a family $\mathcal{R}$ of resolvers. The object of study is
$\mathcal{R} \times \Proc \to$ distributions over normal forms. Two resolvers
are \emph{separated} at $P$ when their images at $P$ differ. Quantum advantage,
if it exists, is a separation, not a full-abstraction theorem.
\end{definition}

\begin{definition}[Support adequacy and the interference
deficit]\label{def:deficit}
The one-sided obligation on a resolver is
$\supp(\sem{P}) \subseteq \Reach(P)$. The \emph{interference deficit}
$\mathcal{D}(P) := \Reach(P) \setminus \supp(\sem{P})$ is the set of
operationally reachable normal forms carrying amplitude zero. Destructive
interference is exactly the failure of the reverse inclusion.
\end{definition}

Contexts remain useful, but as apparatus rather than as observers: an
arrangement that converts a phase difference into a difference of support is an
interferometer, and \S\ref{sec:erasure} builds the smallest one.

%%%%%%%%%%%%%%%%%%%%%%%%%%%%%%%%%%%%%%%%%%%%%%%%%%%%%%%%%%%%%%%%%%%%%%%%%%%%%%%
\section{What unmodified rho already gives}\label{sec:stage1}

Read \eqref{eq:race} as an equational constraint with uniform weights. The
question is which derivations can reconverge, since interference requires
reconvergence and a non-zero relative phase and nothing else.

The previous version asserted that reconvergence in unmodified rho requires
structurally congruent racing messages. That is false, and the counterexample is
instructive. Take $Q_1 \not\scong Q_2$ and
\[
  x!(Q_1) \mid x!(Q_2)
  \mid \textsf{for}(y \leftarrow x)\{a!(0)\}
  \mid \textsf{for}(z \leftarrow x)\{b!(0)\} .
\]
The two complete matchings send $Q_1$ to $y$ and $Q_2$ to $z$, or the reverse.
Both yield $a!(0) \mid b!(0)$. Non-congruent racers reconverge.

They reconverge because $y$ and $z$ are \emph{discarded}. That is the general
mechanism, and it is worth stating in the form that survives.

\begin{proposition}[Stage one, restated]\label{prop:stage1}
Under \eqref{eq:race} with uniform non-negative weights, two derivations
reconverge exactly when the outcome does not record which datum went where.
Congruence of the racing messages is one way to achieve this; discarding the
bound names is another. In either case the interference is constructive, since
non-negative reals cannot cancel.
\end{proposition}

The example is not an anomaly but an instance: it is a \emph{distributed join},
the case of Theorem~\ref{thm:scattering} in which the patterns of a contention
set are spread over several receipts rather than gathered into one. The
constructive half of the old proposition survives; the confinement to congruent
racers does not, and its loss is a gain --- unmodified rho interferes more often
than we had claimed, and always for the same reason.

%%%%%%%%%%%%%%%%%%%%%%%%%%%%%%%%%%%%%%%%%%%%%%%%%%%%%%%%%%%%%%%%%%%%%%%%%%%%%%%
\section{The graded clause, in brief}\label{sec:design}

The design is developed in \cite{ResAlg}; this section fixes notation.

The \texttt{where} clause is two-sorted. The \emph{crisp} sort $\beta$ is the
existing condition language, with the full Boolean structure. The \emph{graded}
sort takes values in a commutative semiring $\V$:
\[
  \psi ::= \crisp{\beta} \mid v \mid \psi \tensorv \psi \mid \psi \sumv \psi
  \mid \langle K \rangle_{\V}\,\psi \qquad (v \in V).
\]
There is no negation on the graded sort, since $\Cx$ has no canonical one, and
$\sumv$ is where interference enters the logic: a disjunction may be false with
both disjuncts true.

\begin{definition}[Candidates and per-match evaluation]\label{def:cand}
At a receipt site $r$ in a term $P$, $\Cand(P,r)$ is the set of injective
assignments of $r$'s patterns to data admissible at its channels. The clause is
evaluated once per candidate, not once per communication, and
$\sem{\psi_r}(\sigma) \in V$ is its value at $\sigma$. A candidate with value
zero does not fire, so the crisp sort keeps its gating job and the graded sort
does nothing else.
\end{definition}

Reading a bound name in the clause does not consume it: guard evaluation is
pure, which the existing condition evaluator already enforces by construction.
This exemption is used throughout \S\ref{sec:gadgets} and is what allows a
clause to depend on both the datum received and the assignment chosen.

%%%%%%%%%%%%%%%%%%%%%%%%%%%%%%%%%%%%%%%%%%%%%%%%%%%%%%%%%%%%%%%%%%%%%%%%%%%%%%%
\section{The quantum namespace}\label{sec:namespace}

Existing quantum process calculi separate classical from quantum data and use a
type system to guarantee that each qubit is owned by a unique process
\cite{GayNagarajanCQP,GayNagarajanTypes}. The rho calculus already has the
distinction, in its introduction forms.

\begin{definition}[The two namespaces]\label{def:namespace}
A name is \emph{quantum} if it is bound by \texttt{new}, or received on a
quantum channel. Every other name --- in particular every quotation $@P$ --- is
\emph{classical}. Quantum names travel only on quantum channels. A clause may
take values outside $\Bool$ only at a receipt whose channels are quantum.
\end{definition}

Sorting is by binding structure rather than by syntactic form, because $@(*x)$
is a quotation denoting the \texttt{new}-bound name $x$; reflection makes the
surface form unreliable and the binder does not.

\begin{proposition}[Ownership from unforgeability]\label{prop:ownership}
A \texttt{new}-bound name cannot be denoted by any process that was not given
it, whereas $@P$ can be constructed by anyone who can write $P$. Hence a datum
on a quantum channel has at most one prospective consumer per firing, and no
process can acquire a reference to a quantum channel except by receiving one.
\end{proposition}

Two consequences, both of which the previous version had to argue for.

\begin{remark}[Contention is not ownership]\label{rem:contention-ownership}
Unique ownership does not forbid several receipts listening on one quantum
channel. A race has exactly one winner per datum, so the race is a competition
\emph{over} ownership and not a violation of it. This is the point on which the
present design departs from the quantum process calculi: they resolve process
non-determinism by scheduling and take probability only from measurement, so the
branching of a race is never coherent. Here it is the only place coherence comes
from.
\end{remark}

\begin{remark}[The herald is classical by construction]\label{rem:herald-class}
The herald channel of \S\ref{sec:dilation} is $h(r) := @r$, the quotation of the
site. By Definition~\ref{def:namespace} it is classical. So the boundary between
the coherent and the classical parts of a run is drawn by the calculus, and
reading a herald is a classical receive --- which is where the instrument acts.
\end{remark}

Freshness in the classical namespace is not lost: quoting a process that
mentions a \texttt{new} name, as in $@(c!(*x))$, is fresh and classical.

%%%%%%%%%%%%%%%%%%%%%%%%%%%%%%%%%%%%%%%%%%%%%%%%%%%%%%%%%%%%%%%%%%%%%%%%%%%%%%%
\section{The denotation is register-indexed}\label{sec:register}

A process term cannot be the carrier of a quantum state, for a reason that shows
up as soon as one asks how two of them combine: there is no operation taking two
quantum states to a quantum state. What there is, is a tensor product of states
on \emph{disjoint} systems. So the denotation must be indexed.

\begin{definition}[Register and state space]\label{def:register}
The \emph{register} of $P$, written $\mathrm{reg}(P)$, is the set of quantum
names at which $P$ can produce or consume. A name occurring only as inert
payload is not in the register. For a register $S$,
\[
  H_S \;:=\; \bigotimes_{x \in S} \mathcal{F}(x),
\]
where $\mathcal{F}(x)$ is the Fock space graded by the occupancy of $x$: the
number of messages resting on that channel. $\sem{P} \in H_{\mathrm{reg}(P)}$.
\end{definition}

\begin{proposition}[Bosonic statistics are forced by
congruence]\label{prop:bose}
Parallel composition is associative and commutative, so a configuration records
how many messages sit on a channel and not which is which. Hence the
$n$-occupancy sector of $\mathcal{F}(x)$ is the $n$-th \emph{symmetric} power,
and the amplitude of a coinciding outcome of an $m$-fold contention is a
permanent rather than a determinant.
\end{proposition}

This settles, upstream of any choice of normalisation, a question the previous
version treated as a fork: whether an $m$-fold degeneracy carries weight $m$ or
$m^2$ is not a matter of taste but a consequence of the structural rules. A
calculus whose parallel composition were not commutative would give a different
answer, and there is no such calculus here.

\begin{proposition}[Composition, and when it forgets]\label{prop:compose}
There is a map $H_S \otimes H_T \to H_{S \cup T}$ adding occupancies, defined
for all $S, T$. It is injective if and only if $S \cap T = \emptyset$. On
overlap it identifies $\ket{2}\otimes\ket{0}$, $\ket{1}\otimes\ket{1}$ and
$\ket{0}\otimes\ket{2}$, forgetting which side supplied which quanta.
\end{proposition}

The forgetting is not a defect. It is indistinguishability, and it is where the
permanent of Theorem~\ref{thm:scattering} comes from. What genuinely fails on
overlapping registers is something else:

\begin{remark}[The interaction term]\label{rem:interaction}
$\sem{P \mid Q}$ is not determined by $\sem{P}$ and $\sem{Q}$ when the registers
overlap, because $P \mid Q$ has candidates neither factor has. The failure of
compositionality is coextensive with the presence of a race
(\S\ref{sec:completion}), and is the same phenomenon as an interaction term in a
joint Hamiltonian. Non-compositionality is also the expected condition in this
area rather than a peculiarity of the present design; see \cite{YingYuFeng} and
\S\ref{sec:related}.
\end{remark}

\begin{remark}[Separation]\label{rem:separation}
A partial tensor defined on disjoint registers is a separation algebra, and the
separating conjunction the generated logic of \cite{ResAlg} already carries is
that tensor. The structure was visible on the logic side before it was visible
on the semantics side.
\end{remark}

Allocation and discard are now visible as the two directions of one thing.
\texttt{new} grows the register, adjoining a factor in its vacuum, which is an
isometry. Discard shrinks it, which is a partial trace. \S\ref{sec:dilation}
dilates the second; \S\ref{sec:fragment} identifies when the first suffices.
%%%%%%%%%%%%%%%%%%%%%%%%%%%%%%%%%%%%%%%%%%%%%%%%%%%%%%%%%%%%%%%%%%%%%%%%%%%%%%%
\section{Contention-set completion}\label{sec:completion}

The one-step operator sums over receipt sites as well as over candidates:
\begin{equation}\label{eq:onestep}
  T(P) \;=\; \sum_{r}\ \sum_{\sigma \in \Cand(P,r)}
    \sem{\psi_r}(\sigma)\;\cdot\;\Out(r,\sigma).
\end{equation}
A condition on one site therefore cannot make $T$ norm-preserving. Worse, the
outer sum conflates two different things: candidates at different sites may be
alternatives, or they may be concurrent, and only the first should be summed.

\begin{definition}[Conflict, components, matchings]\label{def:conflict}
Two candidates \emph{conflict} when they share a datum or share a receipt.
A \emph{contention component} at $P$ is a connected component of that relation.
A \emph{step} of a component is one of its \emph{maximal matchings}: a maximal
set of pairwise non-conflicting candidates.
\end{definition}

The terminology follows event structures, where conflict names events that
cannot both occur \cite{Winskel}. It is the phenomenon the Born rule is here to
decide, and not the failure mode; the failure modes of \S\ref{sec:fragment} are
called \emph{collisions}.

\begin{proposition}[Completion]\label{prop:complete}
Let $C_1,\dots,C_k$ be the contention components at $P$. Then $T$ factorises as
$\bigotimes_i T_{C_i}$, and $T$ preserves norm at $P$ if and only if for each
component
\[
  \sum_{o \,\in\, \Out(C_i)}
  \Bigl| \sum_{\substack{m \,\text{a maximal matching of } C_i\\ \Out(m) = o}}
          w(m) \Bigr|^{2} \;=\; 1 .
\]
\end{proposition}

Three things in that statement do work. Components \emph{tensor}: candidates in
different components are concurrent and both fire, so treating them as
alternatives double-counts. The alternatives within a component are its maximal
\emph{matchings}, not its individual candidates --- this is the RSpace zip, the
bag of data and the bag of patterns admissibly pairable. And the norm is taken
\emph{after} summing over the fibre of $\Out$, since matchings with coinciding
outcomes interfere before they are measured.

Every familiar case is an instance. \eqref{eq:race} is one receipt with two
data. A single receipt with $m$ patterns and $m$ data has $m!$ matchings, which
is Theorem~\ref{thm:scattering}. Two receipts competing for one datum is the
partial-matching case, and its completion condition is literally
\eqref{eq:race}'s. The distributed join of \S\ref{sec:stage1} is two receipts
and two data with a one-element fibre.

Measured, over five configurations and four candidate conventions:

\begin{table}[h]
\centering\small
\begin{tabular}{@{}lrrrr@{}}
\toprule
configuration & site & candidate & matching & \textbf{fibre} \\
\midrule
the race equation \eqref{eq:race}     & $1$ & $1$ & $1$ & $\mathbf{1}$ \\
two receipts, one datum               & $2$ & $1$ & $1$ & $\mathbf{1}$ \\
the distributed join                  & $8$ & $4$ & $2$ & $\mathbf{1}$ \\
the two-fold join                     & $2$ & $2$ & $2$ & $\mathbf{1}$ \\
two independent races                 & $2$ & $1$ & $1$ & $\mathbf{1}$ \\
\bottomrule
\end{tabular}
\caption{Squared norm of the image under four normalisation conventions;
$1$ is preservation. Only the fibre convention survives. Normalising over
candidates gets the right number on independent races for the wrong reason: it
treats concurrency as choice.}\label{tab:conventions}
\end{table}

\begin{remark}[Maximal progress on quantum channels]\label{rem:maxprog}
If a component may decline to fire, then ``$c_1$ then $c_2$'' and ``both at
once'' are distinct histories reaching one configuration, and the interleaving
inflation Proposition~\ref{prop:causal} exists to remove returns by another
door. We therefore take maximal progress \emph{on quantum channels}, leaving
classical channels to interleave as usual. Definition~\ref{def:namespace} is
what makes that expressible. It has a defensible reading: under a Born-rule
resolver an enabled race is resolved, and declining is not among the options.
\end{remark}

\begin{proposition}[Causal histories, not interleavings]\label{prop:causal}
Independent redexes must be counted once. Summing over interleavings instead
inflates the squared norm of a layer of $n$ independent sites by $n!$; measured,
the inflation is exactly $2$, $6$ and $24$ for $n = 2,3,4$ against a causal
squared norm of $1.000000$. Under Definition~\ref{def:conflict} and
Remark~\ref{rem:maxprog} the quotient is not imposed on histories: concurrent
firings are bundled into one matching and independent components tensor, so the
inflation cannot arise.
\end{proposition}

\begin{remark}[Two honest costs]\label{rem:costs}
Completion is \emph{nonlocal}: a clause can be normalised only in the light of
what else is contending, so isometry is a property of a configuration and not of
a clause. And the sum over maximal matchings is permanent-shaped, hence $\#P$
--- which is the physics \cite{AaronsonArkhipov} but bites an implementation, so
what a compiler checks must be a syntactic sufficient condition such as a bound
on component size.
\end{remark}

\begin{conjecture}[Global norm preservation]\label{conj:global}
Proposition~\ref{prop:complete} holds at a configuration. That the induced
operator preserves norm along arbitrary concurrent execution, with components
merging and splitting as data are produced and consumed, is conjectured and not
proved.
\end{conjecture}

%%%%%%%%%%%%%%%%%%%%%%%%%%%%%%%%%%%%%%%%%%%%%%%%%%%%%%%%%%%%%%%%%%%%%%%%%%%%%%%
\section{The join is the recombiner}\label{sec:erasure}

Interference needs reconvergence, and reconvergence needs the outcome to forget
which datum went where. No new term former is required.

\begin{definition}[Symmetric continuation]\label{def:sym}
$P$ is \emph{symmetric} in $y_1,\dots,y_m$ when
\[
  P\{@Q_{\sigma(1)}/y_1,\dots\} \;\scong\; P\{@Q_{\tau(1)}/y_1,\dots\}
\]
for all bijections $\sigma,\tau$ and all payloads.
\end{definition}

Symmetry is cheap to arrange, because parallel composition is commutative:
$z!(*y_1) \mid z!(*y_2)$ is symmetric and $z_1!(*y_1) \mid z_2!(*y_2)$ is not.

\begin{theorem}[Scattering]\label{thm:scattering}
Let $C$ be a contention component with $m$ patterns and $m$ data, and let
$A_{ij}$ be the clause value on the matching assigning datum $j$ to pattern $i$.
If the continuations are symmetric then all $m!$ matchings coincide and the
outcome carries $\Per(A)$. If a clause carries the sign of the permutation, the
outcome carries $\Det(A)$ instead.
\end{theorem}

\begin{corollary}[Which-path, syntactically]\label{cor:whichpath}
No clause produces interference at a component whose continuations distinguish
their bound names. Whether a race interferes is decided by the continuation, not
by the weights.
\end{corollary}

The patterns of $C$ need not belong to one receipt. Theorem~\ref{thm:scattering}
covers the distributed join of \S\ref{sec:stage1}, which is why the
counterexample there is an instance rather than a refutation.

\subsection{The minimal interferometer}\label{ssec:interferometer}

Take $m = 2$ and a clause returning $a$ on the identity assignment and $b$ on
the transposition. With a symmetric continuation there is one outcome with
amplitude $a+b$; with a distinguishing continuation there are two, with
$|a|^2$ and $|b|^2$. For $a = b = 1/\sqrt2$ and relative phase $0$, $\pi$,
$\pi/2$ the coherent weights are $2.000000$, $0.000000$, $1.000000$ against an
incoherent $1.000000$ throughout. Three lines of rholang, and it is
Hong--Ou--Mandel \cite{HOM}: $z$ is one detector, $z_1$ and $z_2$ are two.

%%%%%%%%%%%%%%%%%%%%%%%%%%%%%%%%%%%%%%%%%%%%%%%%%%%%%%%%%%%%%%%%%%%%%%%%%%%%%%%
\section{When is $\sem{P}$ a quantum state?}\label{sec:fragment}

Proposition~\ref{prop:complete} makes each column of $T$ a unit vector. That is
half of $T^{\dagger}T = \Id$. The other half is that distinct columns be
orthogonal, and it fails for reasons that have nothing to do with weights.

\begin{definition}[Collision]\label{def:collision}
Two configurations occurring in one reachable superposition \emph{collide} when
their images share support, so that orthogonality becomes a constraint linking
the weights of different configurations. The collision is \emph{severe} when one
image is contained in the other, since orthogonality then requires forbidding a
transition outright.
\end{definition}

\begin{proposition}[Absorption]\label{prop:absorb}
A configuration with no enabled component is a fixed point of $T$ under the
fictitious identity jump. If another configuration in the same superposition
reduces to it, the two columns are nested and no choice of weights makes them
orthogonal. Termination is not a unitary phenomenon.
\end{proposition}

\begin{definition}[Halt tagging]\label{def:halt}
A configuration with no enabled component does not idle; it becomes
\emph{halted}, tagged with the step at which it halted, and is thereafter
frozen.
\end{definition}

The tag records \emph{when}, never \emph{which}: branches halting together share
it, so it is a coarsening of history in the same sense as the herald of
\S\ref{sec:dilation}, and by Corollary~\ref{cor:whichpath} it writes no
which-path record.

\begin{definition}[Spectator]\label{def:spectator}
A token not consumed by a step is a \emph{spectator}. A spectator resting on a
channel that a firing emits to is indistinguishable from the emitted quantum, by
Proposition~\ref{prop:bose}.
\end{definition}

\begin{definition}[Staging]\label{def:staging}
A term is \emph{staged} when every receipt emits only to quantum channels that
no other receipt emits to and that no spectator occupies --- fresh wires per
stage, which is what a circuit is.
\end{definition}

Measured, over six systems, at one to three tokens and depth three, with
collisions counted over all reachable layers:

\begin{table}[h]
\centering\small
\begin{tabular}{@{}lrrrrrr@{}}
\toprule
& \multicolumn{2}{c}{plain}
& \multicolumn{2}{c}{halt tagging}
& \multicolumn{2}{c}{tagging + staging}\\
\cmidrule(lr){2-3}\cmidrule(lr){4-5}\cmidrule(lr){6-7}
system & all & severe & all & severe & all & severe\\
\midrule
cascade                & $47$ & $\mathbf{20}$ & $29$ & $\mathbf{2}$ & $132$ & $\mathbf{0}$\\
join                   & $16$ & $\mathbf{12}$ & $4$  & $\mathbf{0}$ & $6$   & $\mathbf{0}$\\
beamsplitter, spectator& $6$  & $0$ & $6$ & $0$ & $8$ & $\mathbf{0}$\\
dual rail, two stages  & $\mathbf{0}$ & $\mathbf{0}$ & $\mathbf{0}$ & $\mathbf{0}$ & $\mathbf{0}$ & $\mathbf{0}$\\
\bottomrule
\end{tabular}
\caption{Collisions found by exhaustive search. Of the $32$ severe pairs in the
plain column, $30$ have a normal form on one side; of the $39$ non-severe pairs,
all $39$ involve a spectator. Dual rail with staging has none of either, at
every token count and depth searched.}\label{tab:collisions}
\end{table}

\begin{proposition}[The quantum-safe fragment]\label{prop:safe}
Let a term have a fixed register, conserve tokens, discard nothing, be halt
tagged and staged, and use the dual-rail encoding of
Definition~\ref{def:wire}. Then in every configuration searched the columns of
$T$ have pairwise disjoint support, so orthogonality is automatic and unitarity
reduces to Proposition~\ref{prop:complete}, which is local.
\end{proposition}

This is stated as a proposition about the search, not as a theorem. What is
established is the diagnosis: the two ways a term leaves the fragment are
absorption and spectators, and both are visible in the syntax. What is not
established is that no third way exists.

\begin{remark}[Why the previous version saw no symptom]\label{rem:blind}
Circuit-shaped terms on fresh channels have one receipt per contention
component, so the per-site condition of \cite{GWv4} coincides there with
Proposition~\ref{prop:complete}, and they are staged, so they collide nowhere.
The construction was right; the semantics under it was not. A worked example
drawn from the circuit model cannot exercise this part of the theory.
\end{remark}

%%%%%%%%%%%%%%%%%%%%%%%%%%%%%%%%%%%%%%%%%%%%%%%%%%%%%%%%%%%%%%%%%%%%%%%%%%%%%%%
\section{Gadgets, and why linearity costs nothing}\label{sec:gadgets}

\begin{definition}[Linear binding]\label{def:linear}
At a receipt on quantum channels, every bound name is used exactly once in the
continuation. Occurrences in the clause do not count, guard evaluation being
pure.
\end{definition}

Linearity is what makes a step reversible in the bookkeeping sense: weakening is
the partial trace and contraction is the copy, and Hilbert space is monoidal
rather than cartesian, so a semantics with free discard of amplitude-bearing
data cannot be unitary. The previous version declined this restriction on the
ground that it costs the gate construction. It does, under the encoding that
version used.

\begin{proposition}[Weakening-free steps are block-diagonal]\label{prop:blockdiag}
Suppose the payload multiset of the outcome is the consumed multiset together
with the continuation's literals, and suppose distinct input basis states
present distinct payload multisets. Then no outcome lies in the image of two
distinct inputs, the one-step operator is block-diagonal in the input basis, and
no site realises a basis-changing map.
\end{proposition}

The second hypothesis is about the \emph{encoding}. A qubit encoded as which
datum sits on a channel satisfies it. A qubit encoded as which channel holds a
conserved token does not: every basis state presents the same multiset,
differing only in assignment.

\begin{definition}[Dual-rail encoding]\label{def:wire}
A qubit is a pair of quantum channels, each carrying exactly one of two
designated tokens $\mathsf{one}$ and $\mathsf{zero}$. The basis states are
$q_0!(\mathsf{one}) \mid q_1!(\mathsf{zero})$ and
$q_0!(\mathsf{zero}) \mid q_1!(\mathsf{one})$.
\end{definition}

\begin{lemma}[Arbitrary one-qubit gate, linearly]\label{lem:1q}
Let $w$ be a fresh quantum channel carrying the two output rail names, and
consider
\[
  \textsf{for}(\, y \multimap q_a \ \&\ z \multimap q_b
   \ \&\ u_1 \multimap w \ \&\ u_2 \multimap w \ \textsf{where}\ \psi\,)
  \ \{\; u_1!(*y) \mid u_2!(*z) \;\}.
\]
Every bound name is used exactly once and nothing is discarded. There are two
candidates, the assignments of the output rails to $u_1$ and $u_2$; reflection
makes a received name usable as the channel of a send, so the candidate chooses
where each token lands. Both inputs reach both output configurations, so the
outcomes coincide across inputs, and since $\psi$ may read both $y$ and $u_1$
without consuming them, every $2\times2$ matrix over $\Cx$ is realised.
\end{lemma}

Computed by enumerating the site's candidates: the Hadamard, $X$, $Z$, $T$ and a
$50{:}50$ beamsplitter are each realised with worst entry deviation $0$, as are
$500$ pseudorandom unitaries; the realised matrix is unitary to $2.2\times
10^{-16}$ or exactly. Under the payload encoding the same linearity gives row
supports $\{0,1\}$ and $\{2,3\}$ with empty overlap --- block-diagonal, as
Proposition~\ref{prop:blockdiag} says.

\begin{lemma}[Arbitrary two-qubit gate, linearly]\label{lem:2q}
Four input rails, four output rails and four ancilla patterns give $24$
candidates per input and $16$ reachable input/outcome pairs at uniform
multiplicity $4$, so the coinciding candidates form a permanent-shaped sum
rather than an artefact. CNOT is realised exactly, as are $200$ pseudorandom
product unitaries.
\end{lemma}

\begin{remark}[What changed, and what did not]\label{rem:notprin}
The previous version concluded that a design may have structural unitarity or
circuit-shaped expressiveness but not both. That was an artefact of
Definition~\ref{def:wire}'s predecessor. Proposition~\ref{prop:blockdiag}
survives with its hypothesis made explicit, and its role is now to explain why
the encoding had to change.
\end{remark}

%%%%%%%%%%%%%%%%%%%%%%%%%%%%%%%%%%%%%%%%%%%%%%%%%%%%%%%%%%%%%%%%%%%%%%%%%%%%%%%
\section{A worked example: Shor's algorithm}\label{sec:shor}

Every number below is computed by a simulator that realises each gate
\emph{through the dual-rail site}, enumerating that site's candidates rather
than multiplying gate matrices, and the listings are emitted from the same gate
list the simulator runs.

\subsection{The gadgets}

\begin{lstlisting}
// Every channel below is bound by `new`, hence quantum: unforgeable, and the
// only namespace in which a graded clause may take complex values.  Installer
// channels are quantum too; a persistent receipt consumes each datum once and
// uses its bound names once per firing, so persistence never breaks linearity.
// A qubit is ONE token distributed over TWO rails.

contract gate1( r0In, r1In, r0Out, r1Out, m00, m01, m10, m11 ) = {
  new w in {
      w!(*r0Out) | w!(*r1Out)                 // the two output rail names
    | for( y o- *r0In & z o- *r1In & u1 o- w & u2 o- w
           where [ *y == one and *u1 == *r0Out ] * *m00
               + [ *y == one and *u1 == *r1Out ] * *m10
               + [ *y == zero and *u1 == *r0Out ] * *m01
               + [ *y == zero and *u1 == *r1Out ] * *m11 )
      { u1!(*y) | u2!(*z) }
  }
}

contract hadamard( r0In, r1In, r0Out, r1Out ) = {
  gate1!( *r0In, *r1In, *r0Out, *r1Out,
           0.7071067811865476,  0.7071067811865476,
           0.7071067811865476, -0.7071067811865476 )
}

// Two-qubit gates join four input rails with four ancilla patterns carrying the
// four output rail names.  The clause reads the tokens and the assignment, so
// it determines both the input and the output basis state; entries off the
// intended support return 0 and those candidates do not fire.

// `bit` reads a rail token as 0 or 1 and `rail` writes one; both are ordinary
// rholang and neither touches an amplitude.
contract bit( t ) = { if (*t == one) { 0 } else { 1 } }
contract rail( r0, r1, b ) = {
  if (*b == 0) { *r0!(one) | *r1!(zero) } else { *r0!(zero) | *r1!(one) }
}

contract gate2( a0In, a1In, b0In, b1In, a0Out, a1Out, b0Out, b1Out, m ) = {
  new w in {
      w!(*a0Out) | w!(*a1Out) | w!(*b0Out) | w!(*b1Out)
    | for( ya o- *a0In & za o- *a1In & yb o- *b0In & zb o- *b1In
         & u1 o- w & u2 o- w & u3 o- w & u4 o- w
           where entry( *m, basisIn(*ya, *yb), basisOut(*u1, *u3) ) )
      { u1!(*ya) | u2!(*za) | u3!(*yb) | u4!(*zb) }
  }
}
\end{lstlisting}

\subsection{The instance}

$N = 15$, $a = 7$, four counting and four work qubits, hence sixteen rail pairs.
Each wire carries a stage index, because gadgets compose in series on fresh
channels: \texttt{q3\_5\_0} is the first rail of wire $3$ after five gates. The
staging of Definition~\ref{def:staging} is exactly this discipline, so the term
is in the fragment of Proposition~\ref{prop:safe} by construction.

\begin{lstlisting}
new q0_0_0, q0_0_1, q1_0_0, q1_0_1, q2_0_0, q2_0_1, q3_0_0, q3_0_1,
    q4_0_0, q4_0_1, q5_0_0, q5_0_1, q6_0_0, q6_0_1, q7_0_0, q7_0_1,
    q0_1_0, q0_1_1, q1_1_0, q1_1_1, q2_1_0, q2_1_1, q3_1_0, q3_1_1,
    q0_2_0, q0_2_1, q1_2_0, q1_2_1, q2_2_0, q2_2_1, q3_2_0, q3_2_1,
    q4_1_0, q4_1_1, q5_1_0, q5_1_1, q6_1_0, q6_1_1, q7_1_0, q7_1_1,
    q0_3_0, q0_3_1, q1_3_0, q1_3_1, q0_4_0, q0_4_1, q2_3_0, q2_3_1,
    q0_5_0, q0_5_1, q3_3_0, q3_3_1, q0_6_0, q0_6_1, q1_4_0, q1_4_1,
    q2_4_0, q2_4_1, q1_5_0, q1_5_1, q3_4_0, q3_4_1, q1_6_0, q1_6_1,
    q2_5_0, q2_5_1, q3_5_0, q3_5_1, q2_6_0, q2_6_1, q3_6_0, q3_6_1,
    q0_7_0, q0_7_1, q3_7_0, q3_7_1, q1_7_0, q1_7_1, q2_7_0, q2_7_1,
    out in {

    // 1. the input register, one token per rail pair
    q0_0_0!(one) | q0_0_1!(zero) |
    q1_0_0!(one) | q1_0_1!(zero) |
    q2_0_0!(one) | q2_0_1!(zero) |
    q3_0_0!(one) | q3_0_1!(zero) |
    q4_0_0!(one) | q4_0_1!(zero) |
    q5_0_0!(one) | q5_0_1!(zero) |
    q6_0_0!(one) | q6_0_1!(zero) |
    q7_0_0!(zero) | q7_0_1!(one) |

    hadamard!( *q0_0_0, *q0_0_1, *q0_1_0, *q0_1_1 ) |
    hadamard!( *q1_0_0, *q1_0_1, *q1_1_0, *q1_1_1 ) |
    hadamard!( *q2_0_0, *q2_0_1, *q2_1_0, *q2_1_1 ) |
    hadamard!( *q3_0_0, *q3_0_1, *q3_1_0, *q3_1_1 ) |

    // 2. modular exponentiation: one candidate, clause value 1,
    //    so amplitudes pass through untouched.  Every rail is
    //    consumed and every token re-emitted: no discard.
    for( x0 o- *q0_1_0 & x0' o- *q0_1_1
       & x1 o- *q1_1_0 & x1' o- *q1_1_1
       & x2 o- *q2_1_0 & x2' o- *q2_1_1
       & x3 o- *q3_1_0 & x3' o- *q3_1_1
       & y0 o- *q4_0_0 & y0' o- *q4_0_1
       & y1 o- *q5_0_0 & y1' o- *q5_0_1
       & y2 o- *q6_0_0 & y2' o- *q6_0_1
       & y3 o- *q7_0_0 & y3' o- *q7_0_1 ) {
      new r in {
        modexp!( *r, 7,
                 8 * bit(*x0) + 4 * bit(*x1) + 2 * bit(*x2) + bit(*x3), 15 ) |
        for( m <- r ) {
          q0_2_0!(*x0) | q0_2_1!(*x0') | q1_2_0!(*x1) | q1_2_1!(*x1') |
          q2_2_0!(*x2) | q2_2_1!(*x2') | q3_2_0!(*x3) | q3_2_1!(*x3') |
          rail( q4_1_0, q4_1_1, (*m / 8) % 2 ) |
          rail( q5_1_0, q5_1_1, (*m / 4) % 2 ) |
          rail( q6_1_0, q6_1_1, (*m / 2) % 2 ) |
          rail( q7_1_0, q7_1_1, *m % 2 ) |
          done!(Nil)
        }
      }
    } |

    // 3. the transform on the counting register
    hadamard!( *q0_2_0, *q0_2_1, *q0_3_0, *q0_3_1 ) |
    gate2!( *q1_2_0, *q1_2_1, *q0_3_0, *q0_3_1,
            *q1_3_0, *q1_3_1, *q0_4_0, *q0_4_1, cphase(2*pi/4) ) |
    gate2!( *q2_2_0, *q2_2_1, *q0_4_0, *q0_4_1,
            *q2_3_0, *q2_3_1, *q0_5_0, *q0_5_1, cphase(2*pi/8) ) |
    gate2!( *q3_2_0, *q3_2_1, *q0_5_0, *q0_5_1,
            *q3_3_0, *q3_3_1, *q0_6_0, *q0_6_1, cphase(2*pi/16) ) |
    hadamard!( *q1_3_0, *q1_3_1, *q1_4_0, *q1_4_1 ) |
    gate2!( *q2_3_0, *q2_3_1, *q1_4_0, *q1_4_1,
            *q2_4_0, *q2_4_1, *q1_5_0, *q1_5_1, cphase(2*pi/4) ) |
    gate2!( *q3_3_0, *q3_3_1, *q1_5_0, *q1_5_1,
            *q3_4_0, *q3_4_1, *q1_6_0, *q1_6_1, cphase(2*pi/8) ) |
    hadamard!( *q2_4_0, *q2_4_1, *q2_5_0, *q2_5_1 ) |
    gate2!( *q3_4_0, *q3_4_1, *q2_5_0, *q2_5_1,
            *q3_5_0, *q3_5_1, *q2_6_0, *q2_6_1, cphase(2*pi/4) ) |
    hadamard!( *q3_5_0, *q3_5_1, *q3_6_0, *q3_6_1 ) |
    gate2!( *q0_6_0, *q0_6_1, *q3_6_0, *q3_6_1,
            *q0_7_0, *q0_7_1, *q3_7_0, *q3_7_1, swapMatrix ) |
    gate2!( *q1_6_0, *q1_6_1, *q2_6_0, *q2_6_1,
            *q1_7_0, *q1_7_1, *q2_7_0, *q2_7_1, swapMatrix ) |

    // 4. read the counting register
    for( c0 o- *q0_7_0 & c1 o- *q1_7_0 & c2 o- *q2_7_0 & c3 o- *q3_7_0 ) {
      out!( 8 * bit(*c0) + 4 * bit(*c1) + 2 * bit(*c2) + bit(*c3) )
    }
}
\end{lstlisting}

\subsection{What it computes}

\begin{proposition}[Deterministic communication is amplitude-transparent]\label{prop:free}
A component with a single maximal matching and a crisp clause contributes the
value $1$. Amplitudes are unchanged by deterministic reduction, and the length
of a derivation does not affect its weight.
\end{proposition}

\begin{corollary}\label{cor:modexp}
Modular exponentiation is ordinary rholang, running inside every branch at no
cost in the semantics. In the listing above it consumes every rail and re-emits
every token, so it discards nothing and stays inside
Definition~\ref{def:linear}.
\end{corollary}

The run fires $8$ one-qubit and $8$ two-qubit dual-rail sites. The worst
deviation between the matrix a gadget realises and the matrix intended is $0$.
The counting register reads $0$, $4$, $8$ and $12$ at $0.250000000$ each and
everything else at zero, totalling $1.000000000000$; continued fractions give
$r = 4$, and $\gcd(7^2-1,15) = 3$, $\gcd(7^2+1,15) = 5$. The worst deviation
from the payload-encoded run of \cite{GWv4} is $0$: the encoding changed and the
algorithm did not.

Twelve of the sixteen reachable outcomes carry amplitude zero. By
Definition~\ref{def:deficit} they are the interference deficit, so Shor's
algorithm in this semantics \emph{is} a large deficit arranged on purpose.

\subsection{What the example shows, and what it does not}

It shows that linearity is not the obstacle and that the fragment of
Proposition~\ref{prop:safe} is expressive enough for the standard algorithms.

It does not show that graded rho has a quantum idiom of its own: every gate is a
hand-built gadget and the encoding is circuit-shaped. By
Remark~\ref{rem:blind} a circuit-shaped example also cannot exercise
\S\S\ref{sec:completion}--\ref{sec:fragment}, which is precisely why those
sections are tested by search rather than by this example. The one primitive the
design offers that a circuit does not hand over is
Theorem~\ref{thm:scattering}: an $m$-fold contention is a BosonSampling instance
\cite{AaronsonArkhipov} in a single term former, and since matchings are
permutation sets the ambient symmetry group is $S_m$ rather than $\mathbb{Z}_N$.

Finally, writing Shor is not running Shor. Classical evaluation of the
denotation is exponential; the value of the term is as a compilation source.

%%%%%%%%%%%%%%%%%%%%%%%%%%%%%%%%%%%%%%%%%%%%%%%%%%%%%%%%%%%%%%%%%%%%%%%%%%%%%%%
\section{Dilation, heralds, and the price of erasure}\label{sec:dilation}

Nothing so far forbids a clause whose matrix is not an isometry, and
renormalising at observation is post-selection.

\begin{proposition}[Overshoot]\label{prop:postbqp}
With $\V = \Cx$ and arbitrary clauses, the graded fragment with the Born reading
decides $\mathsf{PostBQP}$, hence $\mathsf{PP}$ \cite{Aaronson}.
\end{proposition}

\begin{definition}[Contractive clause]\label{def:contractive}
A clause is \emph{contractive} at a component when its matrix satisfies
$\|A\| \leq 1$. Contractivity is a well-formedness judgement, checked at
elaboration; for a clause in the normal form of \cite{ResAlg} with constant
scalars the matrix is a finite table and the check is by inspection.
\end{definition}

This is the one thing the design refuses, and it must be refused rather than
repaired: a per-site rescaling is not a global scalar, since branches passing
through different numbers of sites acquire different factors.

\begin{definition}[Defect and heralded step]\label{def:heraldedstep}
For a contraction $A$, put $D_A := \sqrt{\Id - A^{\dagger}A}$ and
$D_{A^{\dagger}} := \sqrt{\Id - AA^{\dagger}}$. The heralded step is
\[
  V_A : \ket{\gamma} \longmapsto
    A\ket{\gamma}\otimes\ket{\mathsf{ok}}
    + D_A\ket{\gamma}\otimes\ket{\mathsf{fail}},
\]
with unitary completion the Julia--Halmos two-defect block
$\begin{psmallmatrix} A & D_{A^{\dagger}} \\ D_A & -A^{\dagger}\end{psmallmatrix}$
\cite{SzNagyFoias}. The two-defect form is used because a matrix arising from a
directed rewrite need not be normal.
\end{definition}

\begin{proposition}[The dilation is faithful]\label{prop:dilation}
$V_A$ is an isometry and the block is unitary for every contraction $A$, and the
$\mathsf{ok}$ block is $A$ unaltered --- so on the branch where nothing was lost
the semantics is unchanged. Computed over the filter clause
$\mathrm{diag}(1,10^{-3})$ and $200$ pseudorandom contractions, the worst
deviation of $U^{\dagger}U$ from the identity is $2.220\times10^{-15}$ and the
worst deviation of the $\mathsf{ok}$ block from $A$ is exactly zero.
\end{proposition}

The dilation is indexed by the contention component, not by the site and not by
a global step; Proposition~\ref{prop:causal} forbids the last of these.

\begin{definition}[Heralded reduction]\label{def:heraldedred}
The communication rule is amended so that a firing at component $r$ emits one
further message: $h(r)!(\mathsf{ok})$ on the $A$ branch and
$h(r)!(\mathsf{fail})$ on the defect branch, where $h(r) := @r$. No term former
is added; by Remark~\ref{rem:herald-class} the herald channel is classical.
\end{definition}

\begin{proposition}[The herald carries the defect bit only]\label{prop:defectbit}
If the herald ranges over $\{\mathsf{ok},\mathsf{fail}\}$ the outcomes of
coinciding matchings remain congruent and Theorem~\ref{thm:scattering} is
unaffected. If it names the matching, every fibre of $\Out$ becomes a singleton
and the coherent sum degenerates to the incoherent one at every component of the
program.
\end{proposition}

Measured on the interferometer of \S\ref{ssec:interferometer}: bare weights
$2.000000$, $0.000000$, $1.000000$; with an $\mathsf{ok}$ herald, identical;
with a matching-naming herald, $1.000000$ throughout.

\begin{corollary}[Uniform heralding is free where it should be]\label{cor:inert}
At a component whose clause is an isometry the defect branch carries amplitude
zero, so the herald reads $\mathsf{ok}$ on every matching and is inert. Heralding
every component therefore costs nothing exactly where the static check applies,
which demotes that check from an admission criterion to a compiler optimisation:
it says which heralds may be elided.
\end{corollary}

\begin{remark}[Neither traced nor projected, but addressed]\label{rem:uncompute}
A herald left unread at a deficient component is a record, and leaving it unread
is a decoherence event --- the correct semantics for a program that discards
information, not a malformedness to be refused. Reading it is the retry, and by
Remark~\ref{rem:herald-class} reading it is a classical receive. That is where
the instrument acts, and it is the only place it acts.
\end{remark}

\begin{proposition}[Post-selection is not a step]\label{prop:nopostsel}
Under Definitions~\ref{def:contractive} and \ref{def:heraldedred} no step loses
norm, so the renormalisation at observation has nothing to do, and
$\mathrm{diag}(1,\epsilon)$ is a heralded step followed by a projection onto
$\mathsf{ok}$. Projection is not a construct of the semantics; it is available
only as a retry loop, whose failures are firings.
\end{proposition}

\begin{proposition}[Expected attempts]\label{prop:attempts}
If a component heralds $\mathsf{ok}$ with probability $p$ and the retry loop
restores its inputs before each attempt, the expected number of firings before
an $\mathsf{ok}$ is $1/p$ and the expected charge is $\kappa/p$ for $\kappa$ the
cost of one firing.
\end{proposition}

Applied to a Hadamard layer on $n$ wires followed by $\mathrm{diag}(1,10^{-3})$
on each: surviving norms $0.250001$, $0.125000$, $0.062500$, $0.031250$ and
success probabilities $0.062500$, $0.015625$, $0.003906$, $0.000977$ for
$n = 4,6,8,10$, hence expected attempts $16$, $64$, $256$, $1024$ --- exactly
$2^{n}$. The filter that was free costs $2^n$ firings.

\begin{conjecture}[Containment by the ledger]\label{conj:ledger}
A graded program funded by a budget polynomial in the size of the term does not
decide $\mathsf{PP}$, since the post-selection witnessing
Proposition~\ref{prop:postbqp} requires an exponentially small success
probability, hence by Proposition~\ref{prop:attempts} an exponentially large
expected charge.
\end{conjecture}

Proposition~\ref{prop:attempts} assumes independent attempts. An adaptive retry
schedule is not geometric, and adaptivity is exactly the resource that makes
linear optics universal \cite{KLM}; the two demands meet somewhere that has not
been located.

%%%%%%%%%%%%%%%%%%%%%%%%%%%%%%%%%%%%%%%%%%%%%%%%%%%%%%%%%%%%%%%%%%%%%%%%%%%%%%%
\section{Related work, and where this sits}\label{sec:related}

The natural comparison class is not quantum process calculi but \emph{quantum
control}.

CQP \cite{GayNagarajanCQP,GayNagarajanTypes} adds measurement and unitary
transformation to the pi calculus, transmits qubits along channels, and uses a
type system to guarantee that each qubit has a unique owner. Its operational
semantics combines non-determinism arising exactly as in the pi calculus with
the probabilistic results of measurement; the quantum state is a separate
component of the configuration and norm preservation is an invariant of that
component. Process branching is never coherent. Later work carries distributions
over configurations \cite{QuantumBarbs}, but they remain probability
distributions. That architecture does not face the question of this paper,
because it declines the premise: here the branching of a race \emph{is} the
quantum state.

The literature that does face it is quantum alternation. QML
\cite{AltenkirchGrattage} superposes branches under an orthogonality judgement;
QGCL \cite{YingYuFeng} does it with an auxiliary system of quantum coins ---
a control register, which is the branch register of the compilation route ---
and notes that the resulting superoperator semantics is not compositional.
Bădescu and Panangaden \cite{BadescuPanangaden} keep the control superposed and
show that the semantics must then be given by Kraus decompositions rather than
superoperators.

Three of their findings bear on this design. Proposition~\ref{prop:noext} is one
of them, and it converts \S\ref{sec:fa}'s negative result from an argument into
a citation. Non-compositionality is the expected condition in this area, which
is what Remark~\ref{rem:interaction} leans on. And their proposition that
quantum alternation is not monotone in the Löwner order, hence incompatible with
the standard fixed-point semantics for recursion, applies here and is not
repaired by anything in this paper; see Obligation~\ref{ob:recursion}.

Their suggested way out is worth quoting in substance. They propose looking at
quantum optics, where a Mach--Zehnder interferometer is a physically real
alternation, and they observe that in such a device the system being split is
the system the alternate operations act on, with no distinct control qubit. That
is the architecture developed here: the race is the split, there is no control
register, and Definition~\ref{def:wire} makes the interferometer literal. Ying's
other route to recursion, via second quantisation \cite{YingSecondQ}, is also
where Proposition~\ref{prop:bose} lands.

The weighted-history pattern itself --- multiply along a path, sum across paths
--- is standard in semiring-weighted computation \cite{DrosteKuichVogler}, and
stochastic process calculi are long established \cite{Priami,SPiM}. The specific
claim here is narrower and, we think, new: a semiring-valued expression placed
exactly at candidate-match resolution, with rho's joins as explicit recombiners
of weighted alternative histories.

%%%%%%%%%%%%%%%%%%%%%%%%%%%%%%%%%%%%%%%%%%%%%%%%%%%%%%%%%%%%%%%%%%%%%%%%%%%%%%%
\section{Open problems}\label{sec:open}

\begin{obligation}[Global norm preservation]\label{ob:global}
Conjecture~\ref{conj:global}: prove Proposition~\ref{prop:complete} along
arbitrary concurrent execution, with components merging and splitting.
\end{obligation}

\begin{obligation}[Staging as a theorem]\label{ob:staging}
Table~\ref{tab:collisions} is a search. Prove that a halt-tagged, staged,
dual-rail term has pairwise disjoint columns, or find the third failure mode.
\end{obligation}

\begin{obligation}[Recursion and the Löwner order]\label{ob:recursion}
Quantum alternation is not monotone in the standard order on completely positive
maps \cite{BadescuPanangaden}, so it is incompatible with the usual fixed-point
semantics for recursion. Rho's recursion is by reflection, and the budget in the
denotation is precisely a stand-in for a fixed point. This is the deepest open
problem here and nothing above addresses it.
\end{obligation}

\begin{obligation}[The price, made exact]\label{ob:price}
Conjecture~\ref{conj:ledger} assumes independent attempts, and the charge for a
herald has not been fixed against the cost model. Until it is,
``polynomially funded'' names an intention.
\end{obligation}

\begin{obligation}[Herald garbage]\label{ob:garbage}
Heralds at deficient components accumulate. Collecting them without reading them
--- which commits to an outcome --- or discarding them --- which is the trace
this design avoids --- is open.
\end{obligation}

\begin{obligation}[Categorical status]\label{ob:cat}
The register-indexed denotation with its partial tensor is a separation algebra
valued in Fock spaces. Making Remark~\ref{rem:separation} precise, and relating
it to the generated logic's separating conjunction, is the natural next step.
\end{obligation}

%%%%%%%%%%%%%%%%%%%%%%%%%%%%%%%%%%%%%%%%%%%%%%%%%%%%%%%%%%%%%%%%%%%%%%%%%%%%%%%
\section{Conclusion}\label{sec:conc}

The rho calculus does not say how its races are settled, and quantum mechanics
is one of the available answers. Taking that seriously turns out to require four
things that were not in the original design.

A term is not a quantum state. The denotation must be indexed by a register of
unforgeable names, valued in Fock spaces over those names, with a tensor that is
partial and forgets which side supplied which quanta when the registers overlap.
The forgetting is indistinguishability; and because parallel composition is
associative-commutative, bosonic statistics --- and the permanent of the
scattering matrix --- follow from the structural rules rather than from a choice.

The unit at which amplitudes must be completed is the contention set: connected
components of the conflict relation, with maximal matchings as the alternatives
and the fibre-summed amplitudes normalised. Components tensor. The race equation
we began from is the smallest instance.

Not every term denotes a state, and the ways of failing are syntactic.
Absorption --- termination, which is not a unitary phenomenon --- and spectators
--- parked quanta indistinguishable from emitted ones --- are the two we have
found, and halt tagging, staging and a dual-rail encoding remove them in every
case we have been able to construct.

And what cannot be excluded can be priced. Contractivity is refused; the
residual defect is dilated at the component; the ancilla is neither traced nor
projected but written, as a message in the classical namespace, which is where
measurement was always going to have to happen. A programmer who wants
post-selection may have it, as a retry loop, at $2^n$ firings.

What remains is not small. Norm preservation is proved at a configuration and
conjectured in general; staging is measured and not proved; and quantum
alternation's incompatibility with the standard semantics of recursion is a
problem this design inherits and does not solve. But the shape of the object is
now clear enough to say what would settle it.

%%%%%%%%%%%%%%%%%%%%%%%%%%%%%%%%%%%%%%%%%%%%%%%%%%%%%%%%%%%%%%%%%%%%%%%%%%%%%%%
\begin{thebibliography}{99}\small

\bibitem{Aaronson} S.~Aaronson.
Quantum computing, postselection, and probabilistic polynomial-time.
\emph{Proc. Roy. Soc. A} 461:3473--3482, 2005.

\bibitem{AaronsonArkhipov} S.~Aaronson and A.~Arkhipov.
The computational complexity of linear optics. \emph{STOC}, 2011.

\bibitem{AltenkirchGrattage} T.~Altenkirch and J.~Grattage.
A functional quantum programming language. \emph{LICS}, 2005.

\bibitem{BadescuPanangaden} C.~B\u{a}descu and P.~Panangaden.
Quantum alternation: prospects and problems.
\emph{QPL 2015}, EPTCS 195:33--42.

\bibitem{Bennett} C.~H.~Bennett.
Logical reversibility of computation.
\emph{IBM J. Res. Develop.} 17(6):525--532, 1973.

\bibitem{CoeckeDuncan} B.~Coecke and R.~Duncan.
Interacting quantum observables: categorical algebra and diagrammatics.
\emph{New J. Phys.} 13:043016, 2011.

\bibitem{DrosteKuichVogler} M.~Droste, W.~Kuich and H.~Vogler, editors.
\emph{Handbook of Weighted Automata}. Springer, 2009.

\bibitem{GayNagarajanCQP} S.~J.~Gay and R.~Nagarajan.
Communicating quantum processes. \emph{POPL}, 145--157, 2005.

\bibitem{GayNagarajanTypes} S.~J.~Gay and R.~Nagarajan.
Types and typechecking for communicating quantum processes.
\emph{Math. Struct. Comp. Sci.} 16(3):375--406, 2006.

\bibitem{Gillespie} D.~T.~Gillespie.
Exact stochastic simulation of coupled chemical reactions.
\emph{J. Phys. Chem.} 81(25):2340--2361, 1977.

\bibitem{Girard} J.-Y.~Girard.
Linear logic. \emph{Theoret. Comput. Sci.} 50(1):1--102, 1987.

\bibitem{Goertzel} B.~Goertzel, M.~Ikl\'e, I.~F.~Goertzel and A.~Heljakka.
\emph{Probabilistic Logic Networks}. Springer, 2009.

\bibitem{GWv4} L.~G.~Meredith.
Graded where-clauses, version 4. F1R3FLY.io working note, 2026.

\bibitem{HOM} C.~K.~Hong, Z.~Y.~Ou and L.~Mandel.
Measurement of subpicosecond time intervals between two photons by
interference. \emph{Phys. Rev. Lett.} 59(18):2044--2046, 1987.

\bibitem{KLM} E.~Knill, R.~Laflamme and G.~J.~Milburn.
A scheme for efficient quantum computation with linear optics.
\emph{Nature} 409:46--52, 2001.

\bibitem{Kraus} K.~Kraus.
\emph{States, Effects and Operations}. Lecture Notes in Physics 190,
Springer, 1983.

\bibitem{Landauer} R.~Landauer.
Irreversibility and heat generation in the computing process.
\emph{IBM J. Res. Develop.} 5(3):183--191, 1961.

\bibitem{MeredithRadestock} L.~G.~Meredith and M.~Radestock.
A reflective higher-order calculus.
\emph{Electron. Notes Theor. Comput. Sci.} 141(5):49--67, 2005.

\bibitem{Milner} R.~Milner, J.~Parrow and D.~Walker.
A calculus of mobile processes, I and II.
\emph{Inform. and Comput.} 100(1):1--77, 1992.

\bibitem{NielsenChuang} M.~A.~Nielsen and I.~L.~Chuang.
\emph{Quantum Computation and Quantum Information}.
Cambridge University Press, 2000.

\bibitem{PLNPaper} L.~G.~Meredith.
Graded where-clauses: the probabilistic-logic instance.
F1R3FLY.io working note, in preparation.

\bibitem{Priami} C.~Priami.
Stochastic pi-calculus. \emph{Computer Journal} 38(7):578--589, 1995.

\bibitem{QuantumBarbs} Quantum bisimilarity via barbs and contexts:
curbing the power of non-deterministic observers. arXiv:2311.06116, 2023.

\bibitem{ResAlg} L.~G.~Meredith.
Resolution algebras: non-crisp truth values as the syntax of non-determinism
resolution in the rho calculus. F1R3FLY.io working note, in preparation.

\bibitem{Reynolds} J.~C.~Reynolds.
Separation logic: a logic for shared mutable data structures.
\emph{LICS}, 55--74, 2002.

\bibitem{Selinger} P.~Selinger.
Towards a quantum programming language.
\emph{Math. Struct. Comp. Sci.} 14(4):527--586, 2004.

\bibitem{SelingerValiron} P.~Selinger and B.~Valiron.
A lambda calculus for quantum computation with classical control.
\emph{TLCA}, 354--368, 2005.

\bibitem{Shor} P.~W.~Shor.
Polynomial-time algorithms for prime factorization and discrete logarithms on a
quantum computer. \emph{SIAM J. Comput.} 26(5):1484--1509, 1997.

\bibitem{SPiM} A.~Phillips and L.~Cardelli.
A correct abstract machine for the stochastic pi-calculus.
\emph{Bioconcur}, 2004.

\bibitem{Stinespring} W.~F.~Stinespring.
Positive functions on $C^{*}$-algebras.
\emph{Proc. Amer. Math. Soc.} 6:211--216, 1955.

\bibitem{SzNagyFoias} B.~Sz.-Nagy and C.~Foias.
\emph{Harmonic Analysis of Operators on Hilbert Space}.
North-Holland, 1970.

\bibitem{Winskel} G.~Winskel.
Event structures. In \emph{Petri Nets: Applications and Relationships to Other
Models of Concurrency}, LNCS 255:325--392, 1987.

\bibitem{YingSecondQ} M.~Ying.
Quantum recursion and second quantisation. arXiv:1405.4443, 2014.

\bibitem{YingYuFeng} M.~Ying, N.~Yu and Y.~Feng.
Alternation in quantum programming: from superposition of data to superposition
of programs. arXiv:1402.5172, 2014.

\end{thebibliography}

\end{document}
